% Formalización reunida el 26 de septiembre de 2026.
% Procedencia y condiciones: COMUN/procedencia_20260926/.
\clearpage
\section{Normalización conjunta de acción, longitud y acoplamiento gravitatorio}
\label{md:rigidez}

La determinación de una magnitud se conserva al componer las lecturas del
mismo estado. Aquí se reúnen el retorno inscrito, la sección de acción, la
velocidad constitutiva y el carácter positivo de ruta. El antecedente común
es la generación APP--TRIT--TPK, con sus hojas, orientación, acarreo, memoria
y prolongaciones; las coordenadas regionales y $\alpha$ intervienen como
salidas de esa generación. La construcción siguiente utiliza esas salidas
antes de efectuar el reconocimiento dimensional de $G$.

El problema de normalización comprende tres operaciones distintas: producir
una longitud desde el retorno registrado, transportar el carácter sin perder
su complemento y comparar las lecturas de energía y longitud que resultan.
Las identidades de esta sección fijan conjuntamente esas operaciones. Sus
pruebas se refieren a fibras finitas o a operadores acotados uniformemente
positivos; el sector nulo conserva su tratamiento separado.

\subsection{Retorno inscrito y sección longitudinal}
\label{md:rigidez:longitud}

En el portador marcado del transporte se conserva el soporte
\[
 \Omega=\{(j,k,q,u,v):j\in\mathbb Z/12\mathbb Z,\ 
 k-j\in\{0,3,6,9\},\ 1\le q\le8,\
 u,v\in\{1,2,4,5,7,8\},\ u<v\}.
\]
El índice $j$ conserva la fase dodecafásica; los cuatro desplazamientos
conservan el cuarto de giro y los restantes índices la incidencia marcada.
Por tanto $N=|\Omega|=12\cdot4\cdot8\binom62=5760$ cuenta posiciones de
incidencia. El espacio de este registro se transporta junto con la memoria
genealógica; su dimensión finita corresponde a este portador particular.

Sean $B:V_U\to V_K$ el inversor de Hadamard transportado,
$B^*B=BB^*=I/4$, $u_K=N^{-1/2}(1,\ldots,1)$,
$P_K=u_Ku_K^*$ y $Q_K=I-P_K$. Definimos
\[
 C_0=Q_KB,\qquad M_0=P_KB,\qquad U=2B.
\]
La ortogonalidad de los dos proyectores demuestra
\[
 C_0^*C_0+M_0^*M_0=I/4,\qquad C_0^*M_0=0,\qquad U^*U=UU^*=I.
\]
Así, $x\mapsto(2C_0x,2M_0x)$ conserva el archivo completo de este
transporte. $U$ es su normalización unitaria; la evolución enriquecida
$U_t$ mantiene su propio dominio y su actualización.

Sea $E_i=e_ie_i^*$ la resolución del registro y
$\Delta_\Omega(A)=\sum_iE_iAE_i$ su expectativa diagonal.

\begin{lemma}[Retorno longitudinal inscrito]
\label{md:rigidez:retorno}
El retorno y su lectura inscrita satisfacen
\[
 C_0C_0^*=Q_K/4,\qquad
 \Delta_\Omega(C_0C_0^*)=r_NI,\qquad
 r_N=\frac{N-1}{4N}=\frac{5759}{23040}.
\]
Con la norma escalar local $n_H=\alpha^{16}$ del conúcleo de $H_4$,
la composición longitudinal sobre una cocadena de longitud $\beta_*$ es
\[
 \mathscr T_L(\beta_*)=
 n_H\bigl(\Delta_\Omega(C_0C_0^*)\otimes I\bigr)\beta_*
 =q_\alpha\beta_*,
 \qquad q_\alpha=\alpha^{16}r_N.
 \label{md:rigidez:qalpha}
\]
\end{lemma}
\begin{proof}
La igualdad $BB^*=I/4$ da $C_0C_0^*=Q_K/4$. Cada entrada diagonal
de $P_K$ vale $1/N$, lo que prueba la expectativa. Multiplicar una
cocadena por ese escalar conserva su tipo longitudinal y da la última
igualdad. Si $\beta_*(\bar a)=-U_a\beta_*(a)$, la imagen satisface la
misma inversión. Asimismo, de
$\beta_m(a)=\sum_jU_j^{-1}\beta_{m+1}(a_j)$ se obtiene la identidad
refinada multiplicando ambos miembros por $q_\alpha$.
\end{proof}

La expectativa posee una realización inscrita explícita:
$We_i=e_i\otimes e_i$ y
$W^*(A\otimes I)W=\Delta_\Omega(A)$. Para
$R_K=Q_K/4$ y $P_W=WW^*$, el complemento
$Z_\perp=(I-P_W)(R_K\otimes I)W$ verifica
\[
 (R_K\otimes I)W=Wr_N+Z_\perp,\qquad W^*Z_\perp=0,\qquad
 Z_\perp^*Z_\perp=\frac{N-1}{16N^2}I.
\]
En efecto, $R_K^2=R_K/4$, de modo que el cuadrado de la norma completa
es $r_N/4$, mientras la componente inscrita tiene cuadrado $r_N^2$.
La diferencia es precisamente el complemento escrito. La conservación de
ese complemento acompaña la lectura longitudinal; sumarlo a otra lectura
cambiaría el observable.

La expectativa es única entre las aplicaciones bimodulares hacia el
álgebra diagonal que fijan sus proyectores: sobre una unidad matricial
$E_{ij}$, la bimodularidad exige
$\mathcal E(E_{ij})=E_i\mathcal E(E_{ij})E_j$, que vale cero si $i\ne j$,
y $\mathcal E(E_{ii})=E_{ii}$. Esta unicidad conserva la resolución marcada.

Con la sección longitudinal de referencia $L_*$ ya especificada, queda
\begin{equation}
 L_\alpha=q_\alpha L_*,\qquad D=L_\alpha U.
 \label{md:rigidez:longitud-def}
\end{equation}
El lector usa el retorno como operador lineal sobre una longitud. Una
lectura cuadrática de la fuente produciría otro tipo de magnitud. La
sección $L_*$, la resolución inscrita y la norma local son antecedentes
explícitos de esta composición.

\subsection{Transporte polar y reciprocidad del mismo carácter}
\label{md:rigidez:polar}

Sea $X$ el carácter operatorio del sector positivo no nulo, conservando la ruta y
sus normalizaciones. En dimensión finita se exige $X>0$; en dimensión
infinita se exige $X$ acotado, autoadjunto y $X\ge\epsilon I$ para algún
$\epsilon>0$. Escribimos $Y=UXU^*$. La representación radial polar es
\[
 R=(DX^2D^*)^{1/2}.
\]
La raíz positiva es única y, al sustituir $D=L_\alpha U$, se obtiene
$R=L_\alpha Y$. Equivalentemente,
$\|Rv\|^2=\|XD^*v\|^2$ para todo $v$: la polarización recupera
$R^2=DX^2D^*$ y la positividad determina $R$. Esta condición especifica
la métrica inducida que se conserva.

La acción $\hbar$ y la velocidad constitutiva $c$ pertenecen a la misma
sección. La realización circular compatible fija
\begin{equation}
 \omega_0=\frac{c}{L_\alpha},\qquad
 t_0=\frac{\pi_{\rm HMT}L_\alpha}{54c},\qquad
 R_{108}=\frac{54}{\pi_{\rm HMT}}ct_0=L_\alpha.
 \label{md:rigidez:reloj}
\end{equation}
La longitud se ha construido antes de expresar el reloj circular. Esta
igualdad permite conservar la notación $R_{108}$ en los desarrollos
posteriores y fija su normalización. Las bases dimensionales se transportan
conjuntamente; $L_*$ y la base de arista $u_L$ tienen funciones diferentes.

\begin{theorem}[Determinación conjunta del acoplamiento radial]
\label{md:rigidez:teorema}
Los lectores sobre el mismo carácter,
\[
 H=\frac{\hbar c}{L_\alpha}Y,\qquad
 \Lambda=L_\alpha Y^{-1},\qquad R=L_\alpha Y,\qquad M=H/c^2,
\]
satisfacen
\begin{equation}
 H\Lambda=\hbar cI,\qquad R\Lambda=L_\alpha^2I,\qquad
 R=\frac{L_\alpha^2}{\hbar c}H.
 \label{md:rigidez:identidades}
\end{equation}
En la identificación dimensional de $R$ como radio gravitatorio reducido,
$R=GM/c^2$, el coeficiente común es único:
\begin{equation}
 \boxed{G=\frac{c^3L_\alpha^2}{\hbar}.}
 \label{md:rigidez:G}
\end{equation}
\end{theorem}
\begin{proof}
La multiplicación de los lectores da las dos primeras identidades.
Multiplicarlas a la derecha por $\Lambda^{-1}$ prueba la tercera.
Sustituir $M=H/c^2$ identifica el coeficiente. Si $G'$ conserva las mismas
lecturas, $(G'-G)M/c^2=0$; en la fibra no nula, la invertibilidad de $M$
implica $G'=G$.
Todos los productos están definidos en la fibra completa bajo las
hipótesis del enunciado.
\end{proof}

La prueba aplica a todos los modos positivos de esa realización, sin elegir
la masa electrónica para normalizar $G$. La identificación del radio
reducido es explícita; el radio $2GM/c^2$ corresponde a otra convención
de lectura de la misma constante. El sector nulo se mantiene fuera del
inverso. La extensión a operadores no acotados requiere un desarrollo de
dominios y límites distinto del teorema presente.

El cuadrado de $L_\alpha$ procede del producto de longitudes conjugadas.
La acción aparece inversamente al eliminar $\Lambda$ entre energía y
longitud. El paso de energía a masa y de masa a radio aporta $c^4$,
que combinado con el $c$ de $H\Lambda=\hbar cI$ produce $c^3$.
Las dos normalizaciones fijan el factor adimensional a uno. El análisis
dimensional verifica los exponentes; la composición determina el coeficiente.

\subsection{Escalas de Planck y especialización del carácter unitario}
\label{md:rigidez:escalas}

La longitud construida, la sección de acción y la velocidad determinan
\[
 t_P=\frac{L_\alpha}{c},\qquad
 m_{P,a}=\frac{\hbar_a}{cL_\alpha},\qquad
 E_{P,a}=\frac{\hbar_a c}{L_\alpha}.
\]
Sustituir \eqref{md:rigidez:G} en
$\sqrt{\hbar_aG_a/c^5}$, $\sqrt{\hbar_ac/G_a}$ y
$\sqrt{\hbar_ac^5/G_a}$ demuestra respectivamente esas identidades;
la positividad fija cada raíz. La longitud antecede al reconocimiento
convencional de estas escalas.
Para un modo de carácter $y>0$,
\[
 m(y)=m_{P,a}y,\qquad r_g(y)=L_\alpha y,\qquad
 \bar\lambda(y)=L_\alpha/y.
\]
La igualdad de las dos longitudes equivale a $y=y^{-1}$ y, por tanto,
a $y=1$. El producto areal se conserva para todos los modos positivos;
la coincidencia de radios caracteriza el modo unitario.

El paso cronológico satisface $t_0=(\pi_{\rm HMT}/54)t_P$;
el período de fase es $108t_0=2\pi_{\rm HMT}t_P$ y
$\omega_{108}=1/t_P$. Un horizonte $R_h=\zeta_hL_\alpha$ conserva
su condición geométrica propia. Para la especialización Schwarzschild
de la misma masa, $R_h=2r_g=2L_\alpha y$.
Al comparar secciones con $L_\alpha,c$ fijos y
$\rho_a=\hbar_b/\hbar_a$, se obtiene
$G_b=G_a/\rho_a$, $m_{P,b}=\rho_a m_{P,a}$ y
$E_{P,b}=\rho_a E_{P,a}$. Esta comparación de secciones conserva
$L_\alpha$ y $t_P$.

\subsection{Área y restitución de los sectores de memoria}
\label{md:rigidez:area}

La incidencia de borde conserva la agrupación de seis trits.
Si $x_j=r_{2j-1}+3r_{2j}$, $r_i\in\{0,1,2\}$, el mapa
\[
 (x_1,x_2,x_3)\longmapsto
 (r_1+3r_2+9r_3,\ r_4+3r_5+9r_6)
\]
es una biyección entre $(\mathbb Z/9\mathbb Z)^3$ y
$(\mathbb Z/27\mathbb Z)^2$. Su inversa extrae los seis dígitos en orden,
conservando el lugar donde se registra el acarreo. En el toro de borde,
el bloque $A=\{0,\ldots,8\}^2$ tiene dos colecciones de enlaces:
\begin{align*}
 \partial A_{90}&=\{((-1,j),(0,j)),((8,j),(9,j)):0\le j\le8\},\\
 \partial A_{120}&=\{((i,-1),(i,0)),((i,8),(i,9)):0\le i\le8\}.
\end{align*}
Cada colección contiene dieciocho enlaces. Sus etiquetas conservan la
procedencia de canal; cada enlace lleva un operador número con espectro
$\{0,1,2\}$. Sean $N_{90},N_{120}$ las sumas correspondientes.

La escala angular areal conserva el contraángulo ya construido:
\[
 \theta_{\rm clk}=\frac{\pi C^*_{\deg}}{1080},\qquad
 a_0=\frac{1+2\cos(\pi/18)}3\,\theta_{\rm clk}.
\]
El prefactor es la traza real normalizada de
$\operatorname{diag}(1,e^{i\pi/18},e^{-i\pi/18})$.
Así $a_0$ conserva el significado de normalización adimensional de área.
La lectura de área se toma en la rama positiva
$\theta_{\rm clk}>0$ y $\gamma>0$.
Con el funcional $\gamma$ de Barbero--Immirzi y la longitud ya determinada,
\[
 \mathcal A=\gamma a_0L_\alpha^2(N_{90}+N_{120}).
\]
Los treinta y seis operadores locales actúan en factores tensoriales
distintos. Por tanto, el espectro de $N=N_{90}+N_{120}$ es
$\{0,\ldots,72\}$ y la multiplicidad de $n$ es el coeficiente de $z^n$
en $(1+z+z^2)^{36}$. La prueba consiste en multiplicar las funciones
generatrices de los factores. La longitud nueva transporta la unidad
areal y conserva esta combinatoria.

\begin{proposition}[Lectura conjunta de área y procedencia sectorial]
La pareja
$N=N_{90}+N_{120}$ y $\mathcal D_{\partial}=90N_{90}+120N_{120}$
restituye los dos sectores:
\[
 N_{90}=4N-\frac{\mathcal D_{\partial}}{30},\qquad
 N_{120}=\frac{\mathcal D_{\partial}}{30}-3N.
\]
\end{proposition}
\begin{proof}
La matriz de la transformación es
$\left(\begin{smallmatrix}1&1\\90&120\end{smallmatrix}\right)$,
con determinante $30$. Su inversa da las fórmulas. La conmutatividad
de ambos operadores número implica la de los dos lectores.
\end{proof}
En la realización modular con
$\rho_A=Z_A^{-1}\exp(-\Delta_{\rm mod}\mathcal D_{\partial})$,
$\Delta_{\rm mod}>0$, el cálculo funcional da
$-\log\rho_A=(\log Z_A)I+\Delta_{\rm mod}\mathcal D_{\partial}$.
Con las normalizaciones conservadas, el par formado por área y
Hamiltoniano modular recupera las dos ocupaciones. La prueba de restitución
vale para cada valor construido de $\Delta_{\rm mod}$; el nuevo lector
longitudinal no redefine ese coeficiente ni el estado modular.
La suma de ocupaciones y la procedencia sectorial expresan observables
distintos. Del mismo modo, los treinta y seis enlaces de respuesta y las
ochenta y una capacidades unitarias del corte cúbico conservan sus funciones.

\subsection{Variación conjunta y reducción compatible de memoria}
\label{md:rigidez:variaciones}

Para colecciones diferenciables en norma y uniformemente positivas en un
entorno, sea $\kappa=L_\alpha^2/(\hbar c)$. Entonces
\[
 \delta R=\delta\kappa\,H+\kappa\,\delta H,\qquad
 \frac{\delta G}{G}=3\frac{\delta c}{c}
 +2\frac{\delta L_\alpha}{L_\alpha}-\frac{\delta\hbar}{\hbar}.
\]
La primera identidad resulta de diferenciar
$R=L_\alpha Y$ y $H=(\hbar c/L_\alpha)Y$; no requiere
$[Y,\delta Y]=0$. Una variación admisible de conexión que actúe en $Y$
y conserve las tres secciones mantiene $\delta G=0$.
La afirmación comprende las variaciones de esta realización radial;
la equivalencia entre funcionales de campo completos tiene sus propios
dominios y condiciones de frontera.

Sea ahora
$Y=\left(\begin{smallmatrix}A&C\\C^*&D_i\end{smallmatrix}\right)>0$
en la separación entre lectura y memoria, con
$S=A-CD_i^{-1}C^*>0$. La extensión minimizante de un vector $b$ es
$(b,-D_i^{-1}C^*b)$, como demuestra completar el cuadrado:
\[
 \left\langle\binom b m,Y\binom b m\right\rangle
 =\langle b,Sb\rangle+
 \langle m+D_i^{-1}C^*b,D_i(m+D_i^{-1}C^*b)\rangle.
\]
Por tanto la reducción coherente da
\[
 H_{\rm eff}=\frac{\hbar c}{L_\alpha}S,\qquad
 R_{\rm eff}=L_\alpha S,\qquad\Lambda_b=L_\alpha S^{-1}.
\]
La compresión de $Y^{-1}$ sobre el bloque de lectura es $S^{-1}$;
la compresión directa de $Y$ es $A$. La memoria distingue esas operaciones.

Si la eliminación dinámica produce una métrica temporal acotada
$Z\ge z_0I$, $z_0>0$, se toma $T=Z^{1/2}$, acotado e invertible.
La normalización canónica exige los transportes duales
\[
 H_c=T^{-*}H_{\rm eff}T^{-1},\qquad
 R_c=T^{-*}R_{\rm eff}T^{-1},\qquad
 \Lambda_c=T\Lambda_bT^*.
\]
Su multiplicación da $H_c\Lambda_c=\hbar cI$ y
$R_c\Lambda_c=L_\alpha^2I$, sin hipótesis de conmutatividad.
En el desarrollo resolvente de una energía con bloques $H_{bb},V,H_{ii}$,
con $H_{ii}$ positivo invertible y
$|\hbar\omega|<\|H_{ii}^{-1}\|^{-1}$,
\[
 K_b(\omega)=H_{bb}-\hbar\omega I
 -V(H_{ii}-\hbar\omega I)^{-1}V^*
 =H_{\rm stat}-\hbar\omega Z+O(\omega^2),
 \qquad Z=I+VH_{ii}^{-2}V^*.
\]
La última expresión es local alrededor de cero y conserva el orden de
aproximación; los productos duales precedentes son identidades exactas
de los lectores normalizados.

La proporcionalidad también se conserva en la eliminación dinámica
completa. Para $z$ tal que $H_{ii}-zI$ sea invertible, definimos
\[
 \mathcal K_H(z)=H_{bb}-zI-V(H_{ii}-zI)^{-1}V^*,\qquad
 \kappa=\frac{L_\alpha^2}{\hbar c}.
\]
La misma descomposición de lectura y memoria aplicada a $R=\kappa H$
produce
\begin{equation}
 \mathcal K_R(\kappa z)=\kappa\mathcal K_H(z).
 \label{md:rigidez:resolvente-exacto}
\end{equation}
En efecto, sus bloques son $R_{bb}=\kappa H_{bb}$,
$R_{bi}=\kappa V$ y $R_{ii}=\kappa H_{ii}$, y
\[
 (R_{ii}-\kappa zI)^{-1}
 =\kappa^{-1}(H_{ii}-zI)^{-1}.
\]
La sustitución en el complemento de Schur demuestra
\eqref{md:rigidez:resolvente-exacto}, conservando el orden de los factores.
Cuando ambos complementos son invertibles, sus inversas satisfacen
$\mathcal K_R(\kappa z)^{-1}=\kappa^{-1}\mathcal K_H(z)^{-1}$.
El parámetro energético es $z=\hbar\omega$ y su imagen en el lector
radial es $\kappa z$: ambos conservan sus dimensiones respectivas.
La identidad vale en todo ese dominio resolvente, por lo que transporta
conjuntamente todos los órdenes de la memoria dinámica. El desarrollo
local anterior describe sus primeros coeficientes; esta igualdad explica
la conservación exacta de la proporción radial al efectuar la reducción.

\subsection{Rigidez temporal y memoria de frontera}
\label{md:rigidez:memoria}

Conservar únicamente el área admite $R\mapsto\zeta R$,
$\Lambda\mapsto\Lambda/\zeta$. Conservar además $H,\hbar,c$ y
$H\Lambda=\hbar cI$ obliga $\zeta=1$. Si también se transforma
$H\mapsto\zeta H$, el cociente $R/H$ y el coeficiente $G$ permanecen
iguales. Las comparaciones deben conservar las mismas observaciones conjuntas.

La fase levantada aporta otra restricción efectiva. En una realización
acotada $H=H_{\rm clk}+D_0^*WD_0$ con $W>0$, la evolución
$P(t)=\exp(-itH/\hbar)$ determina $H=i\hbar\dot P(0)$.
Fijados $H_{\rm clk},D_0,\hbar$ y el reloj, la igualdad de las evoluciones
implica $D_0^*(W'-W)D_0=0$. Para una ruta con extremo inicial fijado,
$D_0$ es triangular invertible y se deduce $W'=W$. Para un grafo general
se determina la forma sobre $\operatorname{Ran}D_0$; extender después el
dominio de residuos modificaría el problema de eliminación.
Un solo retorno $P(\tau)$ admite alias espectrales, mientras la fase
levantada o los tiempos compatibles que tienden a cero los distinguen.

En la coordenatización de capacidad nonádica, sean
$\rho=\log_{10}9$, $\delta=1-\rho$,
$p_j=\lfloor j/\delta\rfloor$ y $v_j=p_j-j$. La capacidad de una ventana es
\[
 C(j,n)=p_{j+n}-p_j-n.
\]
Para $n=35$ sus valores son $729$ o $730$, según el origen; el retorno de
capacidad módulo nueve distingue el primero dentro de esa clase.
La ventana $j=20\to55$ satisface $p_j:437\to1201$ y
$v_j:417\to1146$. Su fase de capacidad vuelve a $3$ y el cociente aumenta
de $46$ a $127$. Las seis capacidades de $20$ y las veintinueve de $21$
suman $729$; la memoria acumulada sigue siendo parte del estado.

Para la primera llegada $T(k)=\lceil k/\rho\rceil-1$, el recuento es
$N(k,C)=T(k+C)-T(k)-C$. Escribir
$\epsilon(k)=\lceil k/\rho\rceil-k/\rho$ da exactamente
\[
 N(k,C)=(\rho^{-1}-1)C+\epsilon(k+C)-\epsilon(k).
\]
Las ventanas de capacidad $729$ con recuentos $34$ y $35$ conservan
términos de frontera diferentes. Incluso $k=21$ y $k=2289$ comparten
$(k\bmod108,T(k)\bmod108)=(21,22)$ y acumulan respectivamente $1$ y
$109$ vacancias. La coincidencia de calendarios finitos conserva menos
información que el estado completo. Estos cálculos interpretan el retorno;
el lector longitudinal de \eqref{md:rigidez:longitud-def} tiene su propio
constructor y no selecciona su escala mediante esas ventanas.

\subsection{Normalización geométrica y evaluación}
\label{md:rigidez:evaluacion}

El bloque central de APP
$B_c=\left(\begin{smallmatrix}7&2\\2&7\end{smallmatrix}\right)$
satisface $B_c^{\mathsf T}g_cB_c=45g_c$,
$g_c=\operatorname{diag}(1,-1)$. Su normalización positiva unimodular es
$B_c/\sqrt{45}$: un factor positivo adicional $a$ daría determinante $a^2$.
La incidencia areal--volumétrica
$T_{\rm av}=\left(\begin{smallmatrix}1&1\\1&2\end{smallmatrix}\right)$
conserva
$G_{\rm av}=\left(\begin{smallmatrix}2&1\\1&-2\end{smallmatrix}\right)$
por multiplicación directa.
En una realización de Minkowski la forma causal y la soldadura dimensional
tienen papeles distintos. La cocadena
$\beta_m=\ell_m\sigma_mW_\square e_\lambda$,
$\ell_m=\ell_0/9^m$, y la relación $\ell_0=ct_0$ conservan la escala
junto con el cono. Alterar la soldadura conservando sólo la firma cambia
esa realización.

La unidad expansiva de Pell con norma uno y traza cuatro satisface
$x+x^{-1}=4$, $x>1$, y por ello $x=2+\sqrt3$.
La órbita $r_n=r_0x^n$ conserva $r_nr_{-n}=r_0^2$.
Cuando $r_0=L_\alpha$, el producto conserva la misma sección areal.
La norma, la forma marcada, el retorno y la sección longitudinal restringen
objetos diferentes; sus funciones se mantienen al componerlos.

\Needspace{10\baselineskip}
En la sección SI declarada, $L_*=1\,\mathrm m$,
$S_*=10^{-34}\,\mathrm{J\,s}$ y $c=299792458\,\mathrm{m\,s^{-1}}$,
la sustitución estructural produce
\[
 L_\alpha=1.616254982625126330\ldots\,10^{-35}\,\mathrm m,\qquad
 \hbar_{\rm ret}=1.054571816915039061\ldots\,10^{-34}\,\mathrm{J\,s},
\]
\begin{equation}
 G_{\rm ret}=6.674299659414022706\ldots\,10^{-11}
 \,\mathrm{m^3\,kg^{-1}\,s^{-2}}.
 \label{md:rigidez:valor}
\end{equation}
El evaluador sustituye las coordenadas regionales y el registro $K$,
reconstruye $\alpha$, el retorno de $H_5$ y $L_\alpha$; la comparación
metrológica se realiza después. La referencia CODATA 2022 es
$6.67430(15)\,10^{-11}$ en esas unidades \cite{codata2022}.
La diferencia es $-3.40586\,10^{-18}$, equivalente a
$-0.00227$ incertidumbres estándar. Las cifras calculadas describen la
evaluación matemática en esta sección, no una incertidumbre experimental.
Las secciones de acción previa y retornada conservan sus etiquetas y
$\hbar_aG_a=c^3L_\alpha^2$.

El significado estructural de $G$ queda así determinado: es el coeficiente
que convierte, en la realización radial común, la lectura energética en
la longitud conjugada preservando simultáneamente acción y área.
La memoria, las marcas de hoja y las bases forman parte de la construcción
que fija ese coeficiente. Su conservación conjunta explica la unicidad
demostrada.
